\begin{abstract}
We give a detailed exposition of the positive lower density of the
predecessors of a fixed Collatz target. A positive target has a predecessor
set of positive lower natural density if and only if it is not divisible
by three. The proof uses inverse arithmetic, nonperiodic orbit points,
and the analytic inverse-path estimates of Mazur. We state the analytic
inputs separately and derive the counting conclusion step by step,
including the multiplicity bound and the passage from residue averages
to counts of distinct integers. We also prove the resulting criterion
for universal convergence. The positive-density result and its
counterexample consequence occur in Mazur's companion paper; they are
not asserted here as new discoveries. A verification appendix records
the corresponding Lean development.
\end{abstract}
\tableofcontents

\section{Statement and notation}
Let $\N=\{0,1,2,\ldots\}$ and $\N_{>0}=\N\setminus\{0\}$. The ordinary
Collatz map is
\begin{equation}\label{eq:T}
 T(n)=\begin{cases}n/2,&2\mid n,\\3n+1,&2\nmid n.\end{cases}
\end{equation}
We write $n\longrightarrow a$ if $T^k(n)=a$ for some $k\in\N$;
the case $k=0$ is allowed. For $a>0$ and integer $X\ge0$, set
\begin{equation}\label{eq:count}
 R_a=\{n>0:n\longrightarrow a\},\qquad P_a(X)=\#\{n\in R_a:n<X\}.
\end{equation}
Throughout, counts are of distinct positive integers and use a strict cutoff.

\begin{theorem}[Target-density characterization]\label{thm:main}
For every $a\in\N_{>0}$,
\begin{equation}\label{eq:main}
 \left(\exists c_a>0\ \exists X_a\in\N\ \forall X\ge X_a:
 P_a(X)\ge c_aX\right)\quad\Longleftrightarrow\quad 3\nmid a.
\end{equation}
\end{theorem}
The positive direction and the consequence for counterexamples are treated
in \cite{mazur-predecessors}. The analytic construction used below is from
the supporting development of \cite{mazur-density}. We give the elementary
arguments and the deductions from its analytic estimates in full. Those
estimates are cited results: their Fourier and inverse-path construction
is not reproved in this paper. This distinction fixes the scope of the
proof and its dependence on earlier work.

For any set $E\subseteq\N_{>0}$ its lower natural density is
\[
 \underline d(E)=\liminf_{X\to\infty}\frac{\#(E\cap[0,X))}{X}.
\]
An eventual bound by $cX$, with $c>0$, implies $\underline d(E)\ge c$.
Conversely, if $d=\underline d(E)>0$, the definition of a lower limit
gives $\#(E\cap[0,X))/X\ge d/2$ for every sufficiently large $X$.
Thus \eqref{eq:main} is precisely the assertion
$\underline d(R_a)>0\iff3\nmid a$.

\section{Inverse arithmetic in every ternary residue class}
For positive $m$, let $v_p(m)$ denote the exponent of the prime $p$ in $m$.
The odd part is $\operatorname{odd}(m)=m/2^{v_2(m)}$. On positive odd
integers the Syracuse map is
\begin{equation}\label{eq:S}
 S(n)=\operatorname{odd}(3n+1).
\end{equation}
If $v=v_2(3n+1)$, then $T^{v+1}(n)=S(n)$. Repeating this observation
shows that Syracuse reachability implies ordinary Collatz reachability.
We use $Q_a(X)$ for the number of positive odd $n<X$ with $S^k(n)=a$
for some $k\ge0$.

\begin{lemma}[An odd child]\label{lem:child}
Every positive odd $a$ not divisible by three has an odd positive $c$
such that $S(c)=a$.
\end{lemma}
\begin{proof}
If $a\equiv1\pmod3$, put $c=(4a-1)/3$. Otherwise $a\equiv2\pmod3$,
and put $c=(2a-1)/3$. In either case the numerator is positive, odd,
and divisible by three; division by the odd number three preserves
oddness. In the two cases respectively,
$3c+1=4a$ and $3c+1=2a$. Since $a$ is odd, removing exactly the displayed
power of two leaves $a$.
\end{proof}

\begin{lemma}[Valuation calculation]\label{lem:valuation}
For $n>0$, $v_3(4^n-1)=1+v_3(n)$.
\end{lemma}
\begin{proof}
Write $n=3^t u$, where $3\nmid u$. The factorization
\[
 4^u-1=3(1+4+\cdots+4^{u-1})
\]
has second factor congruent to $u\not\equiv0\pmod3$, so its valuation
is exactly one. If $z=1+3h$, then
\[
 z^2+z+1=3(1+3h+3h^2),
\]
which has valuation one. Consequently, whenever $z\equiv1\pmod3$,
\[
 v_3(z^3-1)=v_3(z-1)+1.
\]
Starting at $z=4^u$ and cubing $t$ times proves the formula.
\end{proof}

Define $u_j=(4^j-1)/3=1+4+\cdots+4^{j-1}$, with $u_0=0$.
For $j\ge1$, $u_j$ is odd, $u_j\ge j$, and $S(u_j)=1$.
Direct expansion gives
\begin{equation}\label{eq:uadd}
 u_{m+n}=u_m+4^m u_n.
\end{equation}
For $n>0$, Lemma~\ref{lem:valuation} gives
$v_3(u_n)=v_3(n)$; for $n=0$ divisibility by $3^q$ holds on both sides.
Because $4^m$ is a unit modulo $3^q$, \eqref{eq:uadd} yields
\begin{equation}\label{eq:upermutation}
 u_{m+n}\equiv u_m\pmod{3^q}\quad\Longleftrightarrow\quad3^q\mid n.
\end{equation}
If $0\le i<j<3^q$, their difference $j-i$ is not divisible by $3^q$,
so $u_i,u_j$ have different residues. There are $3^q$ choices and $3^q$
residues; hence $j\mapsto u_j$ is a permutation modulo $3^q$ on this range.
Replacing an index $j$ by $j+t3^q$ preserves its residue and makes $u_j$
arbitrarily large. In particular, every residue contains arbitrarily
large odd one-step predecessors of one.

\begin{lemma}[Affine transport]\label{lem:affine}
Let $c$ be as in Lemma~\ref{lem:child}. Every residue class modulo $3^q$
contains arbitrarily large odd $M$ with $S(M)=a$.
\end{lemma}
\begin{proof}
Set $A=3c+1$ and $F_c(x)=Ax+c$. Expanding both sides gives
\[
 3F_c(x)+1=3Ax+3c+1=A(3x+1).
\]
The identity $\operatorname{odd}(rs)=\operatorname{odd}(r)
\operatorname{odd}(s)$ follows by adding their powers of two.
Thus if $S(x)=1$,
\[
 S(F_c(x))=\operatorname{odd}(A)\operatorname{odd}(3x+1)=a\cdot1=a.
\]
Since $c$ is odd and $A$ even, $F_c(x)$ is odd. It also exceeds $x$.
Finally $A\equiv1\pmod3$, so $A$ has an inverse modulo $3^q$.
For a prescribed residue $r$, choose $x=u_j$ in the residue
$A^{-1}(r-c)$, with $x$ larger than the prescribed size bound.
The preceding permutation argument permits this choice. Then $M=F_c(x)$
has every required property.
\end{proof}

\section{Nonperiodicity and unique hitting times}
A point $x$ is \emph{nonperiodic} for $f$ if $f^k(x)\ne x$ for all $k>0$.
This definition does not exclude a point which later enters a cycle.

\begin{lemma}\label{lem:periodic}
If $f(x)=f(y)$ and both points are periodic, then $x=y$.
If a point is nonperiodic, each forward orbit visits it at most once.
\end{lemma}
\begin{proof}
Choose positive periods $p,q$ for $x,y$. Iteration gives
$f^{pq}(x)=x$ and $f^{pq}(y)=y$. Applying $f^{pq-1}$ to their common
image proves $x=y$. For the second assertion, if
$f^d(z)=f^e(z)=x$ with $d<e$, then
$f^{e-d}(x)=f^{e-d}(f^d(z))=f^e(z)=x$, contradicting nonperiodicity.
\end{proof}

\begin{theorem}[Nonperiodic seed supply]\label{thm:supply}
For odd $a>0$ with $3\nmid a$, for every $q,B\in\N$ and every residue
$r\pmod{3^q}$ there exists $M$ such that
\begin{equation}\label{eq:supply}
 M>B,\quad M\text{ odd},\quad S(M)=a,\quad M\equiv r\pmod{3^q},
 \quad M\text{ nonperiodic for }S.
\end{equation}
\end{theorem}
\begin{proof}
Lemma~\ref{lem:affine} gives $x>B$ and then $y>x$ in the prescribed
residue, with $S(x)=S(y)=a$. Since $x\ne y$, Lemma~\ref{lem:periodic}
prevents both from being periodic. Choose the nonperiodic one.
\end{proof}

\section{Finite measures and a precise analytic input}
We now specify the analytic results used from \cite{mazur-density}.
This section separates the cited estimates from the arguments proved
below; a supply of residues by itself is not a density theorem.

\subsection{Reference densities}
Let $\mu_0$ be the point mass on $\mathbb Z/1\mathbb Z$.
Inductively, sample $Y$ with law $\mu_q$ and, independently, a positive
integer $V$ with $\Pr(V=v)=2^{-v}$. Define $\mu_{q+1}$ as the law of
\begin{equation}\label{eq:pmf}
 2^{-V}(3Y+1)\pmod{3^{q+1}}.
\end{equation}
Multiplication by three makes this independent of the representative
of $Y\pmod{3^q}$; $2^{-V}$ denotes a modular inverse.
Set
\[
 \rho_q(y)=\frac23\,3^q\mu_q(y),\qquad
 \langle g\rangle_q=\frac1{3^q}\sum_{y\bmod3^q}g(y).
\]
The definition gives
\begin{equation}\label{eq:rho}
 0\le\rho_q(y)\le\frac23\,3^q,\qquad
 \langle\rho_q\rangle_q=\frac23.
\end{equation}
The averaging measure is uniform on \emph{all} residues, not only the units.
For $k\ge m$ let $\pi_{k,m}$ be reduction modulo $3^m$.
The fine-scale mixing estimate used here is
\begin{equation}\label{eq:mixing}
 \exists C\ge0\ \forall k\ge m\ge1:\quad
 \left\langle\left|\rho_k-\rho_m\circ\pi_{k,m}\right|\right\rangle_k
 \le\frac{C}{m^2}.
\end{equation}
This is a cited analytic estimate, not a consequence of
\eqref{eq:rho}. The supporting development proves an arbitrary fixed
polynomial rate; only exponent two is needed here.

\subsection{Inverse-path data and the estimates used}
An inverse path ending at an odd seed $M$ is a finite list of odd positive
integers $x_0,\ldots,x_d=M$ with $S(x_i)=x_{i+1}$.
Its valuation word is $(v_0,\ldots,v_{d-1})$, where
\[
 3x_i+1=2^{v_i}x_{i+1},\qquad v_i\ge1.
\]
Once $x_0$ and $d$ are given, this word is uniquely determined.
Conversely, a word and its endpoint determine the path by solving
$x_i=(2^{v_i}x_{i+1}-1)/3$ backwards whenever these are odd integers.

The cited construction organizes such paths into finite weighted families.
We state its required properties in reduced mathematical notation for
the singleton-seed case. This is a combined statement of several estimates
in the supporting development, not a claim that one source theorem has
exactly this presentation. The correspondence to source declarations is
recorded in the reproducibility appendix.

\begin{theorem}[Analytic inverse-path inputs, cited]\label{thm:input}
Fix the initial scale $b=32^5$ and the cap $K_n=\lfloor n/100\rfloor$.
The construction of \cite{mazur-density} has the following properties.
\begin{enumerate}
\item For every odd $M\ge16^b$ it gives a nonnegative sequence $A_n(M)$.
There are a number $p>0$ and nonnegative numbers $d_n$, all independent
of $M$, such that $\sum_{n\ge0}d_n<\infty$ and
\begin{equation}\label{eq:increments}
 |A_{n+1}(M)-A_n(M)|\le d_n.
\end{equation}
For each $n$ there are a modulus $3^{q_n}$ and a residue $r_n$ such that
every sufficiently large odd $M\equiv r_n\pmod{3^{q_n}}$ satisfies
\begin{equation}\label{eq:activate}
 A_n(M)\ge\frac23p.
\end{equation}
The sequence $d_n$ is chosen before the physical seed $M$.
\item A generation $n$ has a finite nonempty label set $I_n$, roots
$r_{n,i}>0$, and a common scale $b_n\ge b+2n$ independent of $M$.
For suitable nonnegative integer shift radii $s_n$, put
\[
 L_{n,i}=\frac{4^{b_n}r_{n,i}}{4\,2^{s_n}3^{b_n}},\qquad
 U_{n,i}=2^{2s_n}L_{n,i}.
\]
Call $x>0$ admissible at $n$ if $L_{n,i}<x\le U_{n,i}$ for every $i$.
For $N\ge4\cdot10^9$, if $x>U_{N,i}$ for every $i$, then $x$ is
admissible at some $n\ge N$.
\item For admissible $(n,x)$ there is a nonnegative terminal mark $B_{n,x}$.
If $\ell(M)=\lim_n A_n(M)>0$, then for some $N_0$,
\begin{equation}\label{eq:mark}
 B_{n,x}\ge\ell(M)/2\qquad(n\ge N_0, x\text{ admissible at }n).
\end{equation}
Write $k_n=\lfloor b_n/4\rfloor$. Positivity of this mark implies
$3^{k_n}\le x$.
\item For each admissible $(n,x)$ there is a finite family $\mathcal Z_{n,x}$
of terminal paths to $M$, with source $s(z)$, weight $w(z)\ge0$, and
\begin{equation}\label{eq:budget}
 x\le s(z)<32x,\qquad s(z)w(z)\le M.
\end{equation}
Two members with the same valuation word are identical. With
$F_j(t)=4^jt+(4^j-1)/3$ and $J_n=\lfloor K_n/2\rfloor+1$, one has
\begin{equation}\label{eq:faninput}
 B_{n,x}\le\sum_{z\in\mathcal Z_{n,x}}w(z)
       \sum_{j=0}^{J_n-1}4^{-j}\rho_{k_n}(F_j(s(z))\bmod3^{k_n}).
\end{equation}
\end{enumerate}
\end{theorem}
The weights describe labeled paths; they are not assumed to be counts or
probabilities on distinct sources. In particular, injectivity of the source
map is deliberately absent from the analytic inputs. We prove it below
using the nonperiodicity of $M$.

The principal analytic ingredients behind this statement are finite-residue
averaging, a positive lower bound on the retained probability product,
a summable variation bound uniform in the initial root, compression of
the first valuation, and synchronization of the physical intervals.
Their proofs and the proof of \eqref{eq:mixing} are the cited portion of
the argument. The following sections give all deductions from those inputs.

\section{Choosing one seed with a positive limiting mark}
\begin{lemma}\label{lem:limit}
If $a$ is odd and $3\nmid a$, there is a nonperiodic odd $M$ with
$S(M)=a$ and $\ell(M)>0$.
\end{lemma}
\begin{proof}
Let $D_N=\sum_{n=N}^{\infty}d_n$. Summability gives $D_N\to0$.
For any fixed allowed seed $M$ and $t>N$, telescoping and
\eqref{eq:increments} give
\[
 |A_t(M)-A_N(M)|
 \le\sum_{n=N}^{t-1}|A_{n+1}(M)-A_n(M)|
 \le\sum_{n=N}^{t-1}d_n\le D_N.
\]
Thus $A_n(M)$ is Cauchy, has a real limit $\ell(M)$, and
\begin{equation}\label{eq:tail}
 |\ell(M)-A_N(M)|\le D_N.
\end{equation}
Choose $N$ such that $D_N<p/3$. Now apply
Theorem~\ref{thm:supply} to the residue $r_N$ modulo $3^{q_N}$, using
a size bound large enough for $M\ge16^b$ and \eqref{eq:activate}.
This chooses one nonperiodic seed $M$ with $S(M)=a$ and
$A_N(M)\ge2p/3$. Consequently
\[
 \ell(M)\ge A_N(M)-D_N>\frac{2p}{3}-\frac p3=\frac p3>0.
\]
All later generations use this same $M$. Uniformity of $D_N$ in $M$
is what permits the order of choices $p,d_n$, then $N$, then $M$.
\end{proof}

\section{From path weights to distinct sources}
Fix the seed $M$ from Lemma~\ref{lem:limit}, and an admissible pair $(n,x)$.
Abbreviate $\mathcal Z=\mathcal Z_{n,x}$ and put
\[
 E=\{s(z):z\in\mathcal Z\},\qquad W=\sum_{z\in\mathcal Z}w(z).
\]
\begin{lemma}[Multiplicity and weighted counting]\label{lem:charge}
The map $z\mapsto s(z)$ is injective. For every nonnegative function
$\psi$ on the positive integers,
\begin{equation}\label{eq:charge}
 \sum_{z\in\mathcal Z}w(z)\psi(s(z))
 \le\frac Mx\sum_{s\in E}\psi(s).
\end{equation}
In particular $xW\le M\#E$.
\end{lemma}
\begin{proof}
Suppose two terminal paths have the same source $s$ and lengths $d,e$.
Both hit $M$. Nonperiodicity and Lemma~\ref{lem:periodic} imply $d=e$.
The deterministic Syracuse orbit fixes their valuation words, so the
terminal-path uniqueness in Theorem~\ref{thm:input} makes the members equal.
By \eqref{eq:budget}, $xw(z)\le s(z)w(z)\le M$. Multiply by
$\psi(s(z))\ge0$, sum, and use injectivity:
\[
 x\sum_z w(z)\psi(s(z))\le M\sum_z\psi(s(z))
 =M\sum_{s\in E}\psi(s).
\]
Divide by $x>0$. Taking $\psi=1$ proves the final assertion.
\end{proof}

\begin{lemma}[Residue census]\label{lem:census}
If $Q=3^k\le x$ and $\psi:\mathbb Z/Q\mathbb Z\to\R_{\ge0}$, then
\begin{equation}\label{eq:census}
 \sum_{z\in\mathcal Z}w(z)\psi(s(z)\bmod Q)\le33M\langle\psi\rangle_k.
\end{equation}
\end{lemma}
\begin{proof}
For any residue, the number of nonnegative integers $t<32x$ in that
residue is at most $32x/Q+1\le33x/Q$. Since $E\subset[x,32x)$,
\[
 \sum_{s\in E}\psi(s\bmod Q)
 \le\frac{33x}{Q}\sum_{y\bmod Q}\psi(y)=33x\langle\psi\rangle_k.
\]
Apply \eqref{eq:charge} and cancel $x$.
\end{proof}

\section{Removing the reference-density mark}
The transformation $t\mapsto4t+1$ has iterates
$F_j(t)=4^jt+(4^j-1)/3$. Indeed this holds at $j=0$, and
$4F_j(t)+1=F_{j+1}(t)$. Since $4^j$ is invertible modulo $3^k$,
each $F_j$ permutes those residues and preserves their uniform mean.

Fix integers $1\le m\le k=k_n$ and suppose $B_{n,x}>0$.
Then $3^k\le x$ by Theorem~\ref{thm:input}. Write
\[
 \delta(y)=|\rho_k(y)-\rho_m(\pi_{k,m}y)|,\qquad
 \varepsilon=\frac{C}{m^2},\qquad H_m=\frac23\,3^m.
\]
Equations \eqref{eq:rho} and \eqref{eq:mixing} imply
\[
 \rho_k(y)\le H_m+\delta(y),\qquad
 \langle\delta\rangle_k\le\varepsilon.
\]
Using the residue census with $\psi=\delta\circ F_j$ gives
\[
 \begin{split}
 \sum_z w(z)\rho_k(F_j(s(z)))
 &\le H_m W+\sum_z w(z)\delta(F_j(s(z)))\\
 &\le H_mW+33M\langle\delta\circ F_j\rangle_k\\
 &=H_mW+33M\langle\delta\rangle_k
 \le H_mW+33M\varepsilon.
 \end{split}
\]
Residues modulo $3^k$ are implicit in these evaluations. Multiply by
$4^{-j}$ and sum over $0\le j<J_n$. Since every summand is nonnegative and
$\sum_{j=0}^{J_n-1}4^{-j}\le\sum_{j=0}^\infty4^{-j}=4/3$,
\eqref{eq:faninput} yields the complete estimate
\begin{equation}\label{eq:unmark}
 B_{n,x}\le\frac43(H_mW+33M\varepsilon)
 =\frac89\,3^m W+\frac{44MC}{m^2}.
\end{equation}
Thus the constants $8/9$ and $44$ arise respectively from
$(4/3)(2/3)$ and $(4/3)33$.

Let $\alpha=\ell(M)/2>0$. Choose an integer
\begin{equation}\label{eq:mchoice}
 m>\max\{1,88MC/\alpha\}.
\end{equation}
As $m^2\ge m$, this implies
\[
 \frac{44MC}{m^2}\le\frac{44MC}{m}<\frac\alpha2
\]
when $MC>0$; if $MC=0$ the same desired upper bound is immediate.
For $n\ge2m$, the inequality $b_n\ge b+2n\ge4m$ gives $k_n\ge m$.
For all sufficiently large $n$ the bound \eqref{eq:mark} also holds.
Combining it with \eqref{eq:unmark}, for every admissible $x$ we get
\[
 \alpha\le B_{n,x}\le\frac89\,3^mW+\frac\alpha2,
 \qquad W\ge\eta:=\frac{\alpha}{2(8/9)3^m}>0.
\]
The integer $m$ and the number $\eta$ have now been fixed independently
of both $n$ and $x$.

\section{A lower bound at every sufficiently large cutoff}
Choose $N$ large enough for the preceding mass bound, with
$N\ge4\cdot10^9$. Since $I_N$ is finite, define
\[
 x_* =\sum_{i\in I_N}\max\{0,U_{N,i}\},\qquad
 Y_* =\left\lceil32(\max\{x_*,0\}+1)\right\rceil.
\]
Take any integer $Y\ge Y_*$ and put $x=Y/32$. Then $x>x_*$ and
$x>U_{N,i}$ for every $i$. Interval synchronization gives some $n\ge N$
at which $x$ is admissible. At this generation $W\ge\eta$.
Every source is odd and satisfies $s(z)<32x=Y$. Its terminal path
reaches $M$, and $S(M)=a$, so $E\subseteq\{s<Y:S^k(s)=a\text{ for some }k\}$.
Lemma~\ref{lem:charge} therefore gives
\[
 \frac{Y}{32}\eta=x\eta\le xW\le M\#E\le M Q_a(Y).
\]
Dividing by $M>0$ proves
\begin{equation}\label{eq:odddensity}
 Q_a(Y)\ge\frac{\eta}{32M}Y\qquad(Y\ge Y_*).
\end{equation}
The synchronization step is why this is a bound for every sufficiently
large $Y$, rather than for a selected sequence of cutoffs.

\section{All positive targets and the divisible-by-three exception}
\begin{proof}[Proof of Theorem~\ref{thm:main}]
First let $a>0$ and $3\nmid a$, allowing $a$ to be even.
The formulas in Lemma~\ref{lem:child} still give an odd positive $c$ and
$v\in\{1,2\}$ with $3c+1=2^v a$. If $3\nmid c$, set $d=c$ and $w=v$.
If $3\mid c$, set $d=4c+1$ and $w=v+2$. In the latter case $d$ is odd,
$d\equiv1\pmod3$, and
\[
 3d+1=12c+4=4(3c+1)=2^{v+2}a.
\]
In both cases $d$ is an odd target coprime to three and
$T^{w+1}(d)=a$: the odd step gives $2^w a$ and $w$ halvings give $a$.
Every odd Syracuse predecessor of $d$ is consequently an ordinary
predecessor of $a$. Equation~\eqref{eq:odddensity}, applied to $d$, gives
$P_a(Y)\ge Q_d(Y)\ge c_dY$ for all sufficiently large $Y$.

Now suppose $3\mid a$. We prove by induction on $k$ that
$T^k(n)=a$ implies $n=2^j a$ for some $j\ge0$.
For $k=0$ take $j=0$. For $k+1$, the induction hypothesis applied
to $T(n)$ gives $T(n)=2^j a$, which is divisible by three.
An odd $n$ would give $T(n)=3n+1\equiv1\pmod3$, a contradiction.
Hence $n$ is even and $n=2T(n)=2^{j+1}a$.
Conversely $T^j(2^ja)=a$ by $j$ halvings. Thus
\begin{equation}\label{eq:exception}
 R_a=\{2^ja:j\ge0\},\qquad P_a(2^ka)=k.
\end{equation}
For completeness, if $2^ka\le X<2^{k+1}a$, then
$P_a(X)\le k+1$ and
$P_a(X)/X\le(k+1)/(2^ka)\to0$. Thus $R_a$ has natural density zero.
In particular it cannot satisfy an eventual positive linear lower bound.
This proves both directions.
\end{proof}

\section{The exact consequence for possible counterexamples}
Let $\mathcal B=\{n>0:n\not\longrightarrow1\}$ and
$B(X)=\#(\mathcal B\cap[0,X))$.
\begin{lemma}\label{lem:bad}
If $\mathcal B\ne\varnothing$, there is $a\in\mathcal B$ with
$3\nmid a$, and $R_a\subseteq\mathcal B$.
\end{lemma}
\begin{proof}
Starting from $n\in\mathcal B$, divide out its factors of two and reach
an odd $u$. If $3\nmid u$, take $a=u$; otherwise take $a=3u+1$.
Then $3\nmid a$ and $n\longrightarrow a$. The point $a$ cannot reach one,
as concatenating paths would make $n$ reach one.

Suppose $s\longrightarrow a$ and $s\longrightarrow1$, with hitting
times $d,t$ respectively. If $d\le t$, then $T^{t-d}(a)=1$.
If $t<d$, then $a=T^{d-t}(1)$ lies on the cycle $1,4,2$, each point
of which reaches one. Both possibilities contradict $a\in\mathcal B$.
Thus every $s\in R_a$ lies in $\mathcal B$.
\end{proof}

\begin{theorem}[Counterexample density and a conditional criterion]
If $\mathcal B$ is nonempty, then $B(X)\ge cX$ eventually for some $c>0$.
Moreover universal Collatz convergence is equivalent to
\begin{equation}\label{eq:sparse}
 \forall\varepsilon>0\ \forall X_0\in\N\ \exists X\in\N:\quad
 X\ge X_0,\quad X>0,\quad B(X)<\varepsilon X.
\end{equation}
\end{theorem}
\begin{proof}
Choose $a$ by Lemma~\ref{lem:bad}. Theorem~\ref{thm:main} and inclusion
give $B(X)\ge P_a(X)\ge c_aX$ for all sufficiently large $X$.
If every positive integer converges, $\mathcal B=\varnothing$ and
\eqref{eq:sparse} holds with $X=\max\{X_0,1\}$.
Conversely assume \eqref{eq:sparse}. If a counterexample existed, the
first part would supply $c>0$ and a cutoff $X_0$ with $B(X)\ge cX$
for all $X\ge X_0$. Taking $\varepsilon=c$ in \eqref{eq:sparse}
would give $B(X)<cX$ at such a cutoff, a contradiction.
\end{proof}
No proof of the hypothesis \eqref{eq:sparse} is supplied. The theorem
therefore gives a conditional criterion, not a proof of universal convergence.

\begin{thebibliography}{9}
\bibitem{mazur-density}
L. Mazur, \emph{Explicit Positive-Density Collatz Convergence in
Logarithmic Time}, manuscript v2.1, 6 September 2026, with accompanying
formal source development. ProofAtlas.
\url{https://www.proofatlas.ai/formalizations/positive-density-log-time-collatz/}.
\bibitem{mazur-predecessors}
L. Mazur, \emph{Positive Lower Density of Collatz Predecessors},
manuscript v1.0, 6 September 2026. ProofAtlas.
\url{https://www.proofatlas.ai/formalizations/positive-lower-density-collatz-predecessors/}.
\end{thebibliography}

